\documentclass[11pt]{article}

\usepackage{amsmath,amssymb,amsthm,mathrsfs}
\usepackage{geometry}
\usepackage{hyperref}
\usepackage{enumitem}
\usepackage{bbm}
\usepackage{physics}
\usepackage{bm}

\geometry{margin=1in}
\setlength{\parskip}{0.6em}
\setlength{\parindent}{0em}

\hypersetup{
  colorlinks=true,
  linkcolor=blue,
  citecolor=blue,
  urlcolor=blue
}

\title{Multiplicity Engrams as Prime-Graded Subsystem Codes:\\
A Formal Operator-Algebraic Framework and Testable Architecture}

\author{Citizen Gardens \\ \\ The Foundation of Multiplicity}

\date{April 3, 2026}

\begin{document}

\maketitle

\begin{abstract}
We develop a fully operator-algebraic formulation of quantum memory engrams as
prime-graded subsystem codes, unifying the \emph{Quantum Memory Engrams} paper
(prime-encoded eigenvalues, tensor networks, feedback) with the broader
Multiplicity Theory stack (multiplicity functor, prime-number operators,
bounded-block meta-relativity, and tensor-network constraints). Engrams are
defined as metastable density operators supported on a narrow band of joint
eigenspaces of commuting prime-indexed charges. We specify a Lindblad
dynamics with sector dephasing and multiplicity-biased repair, give a
graph-theoretic notion of code distance on the exponent lattice, and formulate
a non-reducibility statement: the arithmetic prime structure realizes a
canonical free-commutative-monoid action on multiplicity sectors and is not
mere relabeling. Finally, we outline a minimal simulation and experimental
protocol (three matched baselines) to falsify or corroborate the claimed
advantage of multiplicity-graded codes.
\end{abstract}
\newpage
\tableofcontents

\section{Executive Summary}

The goal of this report is to crystallize an operator-algebraic, testable
formulation of \emph{quantum memory engrams} as \emph{multiplicity-graded
subsystem codes}.

Building strictly on existing ingredients from the Multiplicity Theory stack
and the paper \emph{Quantum Memory Engrams}:
\begin{itemize}[leftmargin=*]
  \item prime-encoded eigenvalues and eigenvectors;
  \item recursive feedback $\Xi(t)$ and dynamic multiplicity equations;
  \item Lindblad-style dissipation and contraction semigroups;
  \item multiplicity functor $M$ and prime-indexed number operators;
  \item bounded self-adjoint blocks (Meta-Relativity);
  \item tensor networks and multi-scale constraints;
\end{itemize}
we:

\begin{enumerate}[leftmargin=*]
  \item Define a \textbf{multiplicity engram code} as a density operator
        supported on a narrow band of joint eigenspaces of commuting prime
        charges $\{Q_p\}_{p\in P_{\text{active}}}$.
  \item Specify a \textbf{Lindblad master equation} combining system dynamics,
        noise, sector dephasing, and multiplicity-biased repair.
  \item Introduce a \textbf{graph-theoretic code distance} on the exponent
        lattice, induced by elementary leakage generators.
  \item Formulate a \textbf{non-reducibility} statement: admissible
        relabelings must preserve the free-commutative-monoid action, with
        the prime factorization realization providing the canonical faithful
        embedding.
  \item Propose a \textbf{three-baseline simulation/experiment}: (i) the
        multiplicity code, (ii) a generic lattice-labeled code with identical
        adjacency and control budget, and (iii) a standard stabilizer-like
        code under the same leakage model.
\end{enumerate}

The result is not merely a new narrative: it is a falsifiable subsystem-code
architecture whose distinctive feature is that its protection and distance
emerge from prime-based multiplicity structure rather than arbitrary labels.

\section{Background: Quantum Memory Engrams and Multiplicity}

\subsection{Quantum Memory Engrams}

The attached paper \emph{Quantum Memory Engrams} models a memory engram as a
time-evolving quantum state
\begin{equation}
  \Psi(t)
  = \sum_{i=1}^N \lambda_i(t) v_i(t) + \xi(t),
  \label{eq:psi-decomposition}
\end{equation}
where:
\begin{itemize}[leftmargin=*]
  \item $\lambda_i(t)$ are time-dependent eigenvalues, encoded via prime
        numbers and evolving with phase coherence;
  \item $v_i(t)$ are eigenvectors representing quantum memory states;
  \item $\xi(t)$ is a stochastic noise term, minimized through feedback.
\end{itemize}
Prime-based encoding is expressed through maps such as
\begin{equation}
  \phi(v_i) = p_k, \qquad
  \phi(e_i) = \prod_{j \in e_i} p_j,
\end{equation}
where $p_k$ are primes, vertices $v_i$ are assigned prime labels, and hyperedges
$e_i$ carry products of primes.

Feedback and stability are introduced via:
\begin{align}
  M(t+1) &= f(M(t), R(t)), \\
  \lambda_i(t+1) &= \lambda_i(t) + \beta \,\Delta\lambda_i(t),
\end{align}
and decoherence/noise correction is expressed at the component level as:
\begin{equation}
  \partial_t \rho_k
  = \alpha_k \rho_k + \gamma_k \sum_j T_{kj} \rho_j
    - \xi_k(t).
\end{equation}
Tensor networks are used to capture interactions:
\begin{equation}
  T(t) = \sum_{i,j,k} T_{ijk} \otimes \phi(p_{ijk}),
\end{equation}
and the eigenvalues exhibit phase coherence
\begin{equation}
  \lambda_i(t) = \lambda_i(0) e^{i \theta_i(t)},
  \qquad
  \theta_i(t) = \omega_i t + \theta_i^0.
\end{equation}

\subsection{Multiplicity Theory Ingredients}

Across the broader Multiplicity Theory stack (FTSA, ZSD, DRMM, Meta-Relativity,
AZ-TFTC, etc.), the following elements are already present:
\begin{itemize}[leftmargin=*]
  \item A multiplicity functor $M$ sending structures to prime factorization
        data, e.g.\ $m = \prod_p p^{\alpha_p}$.
  \item Prime-indexed number operators (or ``multiplicity charges''),
        typically bounded and self-adjoint on block-decomposed spaces.
  \item Lindblad-style dissipative dynamics that preserve positivity and
        generate contraction semigroups.
  \item Tensor-network representations with constraints that encode
        combinatorial or holographic structure.
\end{itemize}

What was missing in the engram context was an explicit, operator-algebraic
connection between:
\begin{enumerate}[leftmargin=*]
  \item Prime encodings at the level of eigenvalues/vectors and tensor
        networks.
  \item A formal notion of code space, distance, and error model that could
        be implemented and falsified on real hardware.
\end{enumerate}

\section{Multiplicity Engram Code: Formal Definition}

\subsection{Prime-Graded Hilbert Space}

Fix a finite set of active primes $P_{\text{active}}$ (consistent with finite
prime truncations used in FTSA/ZSD). For each $p \in P_{\text{active}}$, let
$Q_p$ be a bounded self-adjoint operator on the ambient Hilbert space
$\mathcal H$:
\begin{equation}
  \mathcal Q = \{ Q_p \}_{p \in P_{\text{active}}}, \qquad
  [Q_p, Q_q] = 0, \quad Q_p^\dagger = Q_p.
\end{equation}
The joint spectral decomposition yields a grading
\begin{equation}
  \mathcal H = \bigoplus_{\boldsymbol\alpha} \mathcal H_{\boldsymbol\alpha},
  \qquad
  Q_p \ket{\psi} = \alpha_p \ket{\psi} \quad
  \forall p \in P_{\text{active}},
  \label{eq:joint-eigenspaces}
\end{equation}
where $\boldsymbol\alpha = (\alpha_p)_{p \in P_{\text{active}}} \in
\mathbb N^{|P_{\text{active}}|}$ and $\mathcal H_{\boldsymbol\alpha}$ is the
joint eigenspace with multiplicity exponents $\alpha_p$.

The multiplicity functor $M$ induces a canonical coordinate:
\begin{equation}
  m(\boldsymbol\alpha)
  = \prod_{p \in P_{\text{active}}} p^{\alpha_p},
\end{equation}
realizing the usual Fundamental Theorem of Arithmetic factorization in the
restricted prime set.

\subsection{Engram Support Band}

A \emph{band} of multiplicity vectors is a finite subset
$\mathcal B \subset \mathbb N^{|P_{\text{active}}|}$, chosen according to
the combinatorial structure of the encoded information (e.g.\ central
multiplicity plus a small neighborhood).

\textbf{Definition (Engram support).} The engram support subspace is
\begin{equation}
  \mathcal H_E
  = \bigoplus_{\boldsymbol\alpha \in \mathcal B}
    \mathcal H_{\boldsymbol\alpha}.
\end{equation}

\textbf{Definition (Multiplicity Engram Code).} A \emph{quantum memory engram}
is a metastable density operator $\rho_E$ supported on $\mathcal H_E$ (or
overwhelmingly concentrated there under the dynamics), i.e.
\begin{equation}
  \rho_E \approx \Pi_E \rho_E \Pi_E, \qquad
  \Pi_E = \sum_{\boldsymbol\alpha \in \mathcal B}
          \Pi_{\boldsymbol\alpha},
\end{equation}
where $\Pi_{\boldsymbol\alpha}$ is the projector onto
$\mathcal H_{\boldsymbol\alpha}$.

This formalizes the intuitive notion from the engram paper that a memory
state is a stable pattern in a prime-encoded decomposition, now expressed as
a supported band of multiplicity sectors.

\section{Dynamics: Lindblad Form With Multiplicity Feedback}

We now write the engram dynamics as a Lindblad master equation:
\begin{equation}
  \frac{d\rho}{dt}
  = -i [H_{\text{sys}} + \Xi(t), \rho]
    + \mathcal L_{\text{noise}}(\rho)
    + \mathcal L_{\text{dephase}}(\rho)
    + \mathcal L_{\text{repair}}(\rho).
  \label{eq:master-equation}
\end{equation}
Each term is motivated by, and consistent with, the structures in the engram
paper and related Multiplicity Theory work.

\subsection{System and Feedback}

\begin{itemize}[leftmargin=*]
  \item $H_{\text{sys}}$ is the system Hamiltonian encoding baseline dynamics
        of the memory degrees of freedom.
  \item $\Xi(t)$ is a feedback term, capturing the recursive updates
        $M(t+1) = f(M(t), R(t))$ and spectral corrections
        $\lambda_i(t+1) = \lambda_i(t) + \beta \,\Delta\lambda_i(t)$ in an
        effective Hamiltonian or control term.
\end{itemize}

\subsection{Noise Decomposition}

The total noise is decomposed as
\begin{equation}
  \mathcal L_{\text{noise}}
  = \mathcal L_{\parallel} + \mathcal L_{\perp},
\end{equation}
where:
\begin{itemize}[leftmargin=*]
  \item $\mathcal L_{\parallel}$ commutes with all $Q_p$, thus preserving
        multiplicity sectors.
  \item $\mathcal L_{\perp}$ consists of leakage channels that do not commute
        with the $Q_p$ and thus change the exponent vector
        $\boldsymbol\alpha$.
\end{itemize}

\subsection{Sector Dephasing}

To suppress coherence between distinct multiplicity sectors, we introduce a
dephasing channel:
\begin{equation}
  \mathcal L_{\text{dephase}}(\rho)
  = \sum_{p \in P_{\text{active}}}
      \gamma_p [Q_p, [Q_p, \rho]],
  \qquad \gamma_p > 0.
\end{equation}
This Lindblad form is standard for dephasing in the eigenbasis of
$Q_p$. It damps off-diagonal blocks between different eigenvalues of $Q_p$,
thereby favoring block-diagonal structure in the joint eigenspace
decomposition.

\subsection{Elementary Leakage and Repair Channels}

Let $\boldsymbol e_p$ denote the unit vector in the direction of prime $p$.
We define elementary leakage (and repair) maps that shift the exponent vector:
\begin{equation}
  L_p^{(+)}:\ \mathcal H_{\boldsymbol\alpha}
  \to \mathcal H_{\boldsymbol\alpha + \boldsymbol e_p}, \quad
  L_p^{(-)}:\ \mathcal H_{\boldsymbol\alpha}
  \to \mathcal H_{\boldsymbol\alpha - \boldsymbol e_p},
\end{equation}
whenever the target exponents remain in the truncated lattice (e.g.\ 
$\alpha_p \ge 0$ and within a finite cutoff). The concrete operators can be
specified in a chosen basis; abstractly they generate the free commutative
monoid action on the exponent lattice.

A biased repair channel is then:
\begin{equation}
  \mathcal L_{\text{repair}}(\rho)
  = \sum_{p \in P_{\text{active}}}
    \kappa_p \left( L_p \,\rho\, L_p^\dagger
      - \frac{1}{2} \{L_p^\dagger L_p, \rho\} \right),
  \qquad \kappa_p > 0,
\end{equation}
where $L_p$ collectively denotes a set of $L_p^{(\pm)}$ with bias chosen so
that transitions into $\mathcal B$ are statistically favored. This provides
the operator-algebraic realization of multiplicity feedback and recovery.

By construction, the full generator in \eqref{eq:master-equation} is
completely positive and trace-preserving (CPTP), assuming the usual
Lindblad assumptions.

\section{Logical Structure and Subsystem-Code View}

\subsection{Engram Subspace and Leakage Complement}

We decompose the ambient Hilbert space as
\begin{equation}
  \mathcal H = \mathcal H_E \oplus \mathcal H_E^\perp,
\end{equation}
where $\mathcal H_E$ is the engram support band and
$\mathcal H_E^\perp$ contains all leakage sectors. The projector
$\Pi_E$ defines the engram readout:
\begin{equation}
  \rho_E^{\text{retrieved}} = \Pi_E \rho \Pi_E.
\end{equation}
This is the formal version of the retrieval operation: project onto the
intended multiplicity band and decode within that subspace.

\subsection{Gauge vs Logical Degrees of Freedom}

Within $\mathcal H_E$ one typically expects degeneracy beyond the multiplicity
labels. The \emph{gauge algebra} is generated by operators that:
\begin{itemize}[leftmargin=*]
  \item act within each joint eigenspace $\mathcal H_{\boldsymbol\alpha}$;
  \item preserve the multiplicity labels $\boldsymbol\alpha$.
\end{itemize}
The \emph{logical algebra} is the commutant of the gauge algebra restricted
to $\mathcal H_E$, capturing logical degrees of freedom that are invisible to
the multiplicity charges and gauge actions.

This is the standard subsystem-code perspective: $\mathcal H_E$ factorizes
as a tensor product of a logical subsystem and a gauge subsystem (up to
unitary equivalence), and multiplicity preservation acts as a symmetry
protecting the logical subsystem.

\section{Distance on the Exponent Lattice}

\subsection{Elementary Leakage Graph}

Define the \emph{elementary leakage graph} $G$ whose vertices are exponent
vectors $\boldsymbol\alpha \in \mathbb N^{|P_{\text{active}}|}$, with edges
given by the action of $L_p^{(\pm)}$:
\begin{equation}
  \boldsymbol\alpha \sim \boldsymbol\beta
  \quad \Leftrightarrow \quad
  \boldsymbol\beta = \boldsymbol\alpha \pm \boldsymbol e_p
  \text{ for some } p \in P_{\text{active}}.
\end{equation}
Thus $G$ is an induced subgraph of the integer lattice determined by
finite prime truncation and cutoffs on exponents.

\subsection{Code Distance}

Let $\mathcal B \subset \mathbb N^{|P_{\text{active}}|}$ be the engram band,
and fix a decoding rule that attempts to recover any state supported within
a radius-$R$ neighborhood of $\mathcal B$ in $G$.

\textbf{Definition (Distance).} The code distance $d$ of the multiplicity
engram code is the minimal path length in $G$ required to map a state in
$\mathcal H_E$ to a state that is \emph{logically distinguishable} and lies
outside the correction radius. Equivalently, $d$ is the minimal number of
elementary leakage steps $L_p^{(\pm)}$ such that no allowed recovery map
returns the resulting state to an equivalent state in $\mathcal H_E$.

This definition ties distance to:
\begin{itemize}[leftmargin=*]
  \item the exponent-lattice geometry;
  \item the allowed elementary error set (the $L_p^{(\pm)}$);
  \item the specific decoder used.
\end{itemize}
It generalizes, but does not collapse to, a pure Hamming-distance claim on
the exponents.

\section{Non-Reducibility and Arithmetic Content}

\subsection{Admissible Relabelings}

Consider a map $\sigma: \mathbb N^{|P_{\text{active}}|} \to
\mathbb N^{|P_{\text{active}}|}$ that relabels exponent vectors. We call
$\sigma$ \emph{physically admissible} if:
\begin{enumerate}[leftmargin=*]
  \item it preserves adjacency in the leakage graph $G$, i.e.\ it intertwines
        the action of elementary steps $\boldsymbol e_p$;
  \item it preserves the monoid structure generated by the $L_p^{(\pm)}$, so
        that composition of errors maps coherently under $\sigma$.
\end{enumerate}
This reflects the requirement that the physical noise and repair channels
remain local in the sense of the allowed elementary operations.

\subsection{Prime Factorization as Canonical Realization}

The exponent lattice $(\mathbb N^{|P_{\text{active}}|}, +)$ realizes the free
commutative monoid generated by unit vectors $\boldsymbol e_p$. The
multiplicity map
\begin{equation}
  \boldsymbol\alpha \mapsto m(\boldsymbol\alpha)
  = \prod_{p \in P_{\text{active}}} p^{\alpha_p}
\end{equation}
gives a faithful embedding into $(\mathbb N, \times)$, with:
\begin{equation}
  m(\boldsymbol\alpha + \boldsymbol\beta)
  = m(\boldsymbol\alpha)\,m(\boldsymbol\beta).
\end{equation}
The Fundamental Theorem of Arithmetic guarantees:
\begin{itemize}[leftmargin=*]
  \item uniqueness of factorization;
  \item injectivity of the map $\boldsymbol\alpha \mapsto m(\boldsymbol\alpha)$
        within the truncated prime set.
\end{itemize}

Thus the prime-realized multiplicity $m$ is the canonical coordinate system
for this free commutative monoid. Any admissible relabeling that preserves
the monoid action is equivalent, up to isomorphism, to a multiplicity-based
representation. This shows that the arithmetic structure is not merely a
cosmetic choice of labels: it encodes the semigroup structure of errors and
repairs, and makes composition of transitions additive in exponent space.

\section{Simulation and Experimental Protocol}

\subsection{Three Matched Baselines}

To falsify or corroborate the claimed advantage of multiplicity-graded codes,
we propose a small-dimensional simulation with three baselines:

\begin{enumerate}[leftmargin=*]
  \item \textbf{Multiplicity code:} charges $Q_p$, band $\mathcal B$,
        dephasing $\mathcal L_{\text{dephase}}$, and biased repair
        $\mathcal L_{\text{repair}}$ as defined above.
  \item \textbf{Generic lattice-labeled code:} same Hilbert-space dimension,
        same leakage graph adjacency (same elementary error set), same
        number and strength of control channels, but without the explicit
        prime/multiplicity interpretation.
  \item \textbf{Standard stabilizer-like code:} a stabilizer or decoherence-free
        subspace architecture with comparable locality and resource budget
        under the same leakage model.
\end{enumerate}

The observable is the retained population in the engram band
$\Pi_E \rho(t) \Pi_E$ at a fixed time $t$ (and, optionally, logical fidelity
measures derived from specific logical codewords).

\subsection{Minimal Finite-Dimensional Example}

As a concrete finite example, consider two active primes (2 and 3) and a
$3\times 3$ exponent grid $(\alpha_2, \alpha_3) \in \{0,1,2\}^2$. The Hilbert
space dimension is $d = 9$, with basis states $\ket{\alpha_2,\alpha_3}$.

The charges $Q_2$, $Q_3$ are diagonal in this basis:
\begin{equation}
  Q_2 \ket{\alpha_2,\alpha_3} = \alpha_2 \ket{\alpha_2,\alpha_3}, \quad
  Q_3 \ket{\alpha_2,\alpha_3} = \alpha_3 \ket{\alpha_2,\alpha_3}.
\end{equation}
Choose a simple band, e.g.\ $\mathcal B=\{(1,1)\}$ or a small neighborhood
around it, and define $\Pi_E$ accordingly.

The elementary leakage operators $L_2^{(\pm)}$, $L_3^{(\pm)}$ shift the grid
coordinates by $\pm 1$ along each axis, respecting boundaries.

\subsection{Example Python / QuTiP Skeleton}

Below is a compact code skeleton illustrating how this 9-dimensional system
can be realized in Python with QuTiP. It sets up the basis, charges, and
elementary operators; from there, one can define collapse operators for the
three baselines and run the master equation.

\begin{verbatim}
import qutip as qt
import numpy as np

# 2 primes → 3-level truncation per charge (alpha = 0,1,2)
dim = 3 * 3  # total Hilbert space
# Basis ordering: |a2,a3> with a2,a3 in {0,1,2}
basis = [qt.basis(dim, i*3 + j) for i in range(3) for j in range(3)]

# Multiplicity charges (diagonal in this basis)
Q2 = qt.Qobj(np.diag(np.array([i for i in range(3) for _ in range(3)])))
Q3 = qt.Qobj(np.diag(np.array([j for _ in range(3) for j in range(3)])))

# Engram band B = {(1,1)} (central sector)
engram_states = [basis[1*3 + 1]]  # index 4
Pi_E = sum([psi * psi.dag() for psi in engram_states])

# Elementary leakage operators (weight-1)
def L_p(p_idx, sign):
    # p_idx = 0 for Q2, 1 for Q3; sign = +1 or -1
    mat = np.zeros((dim, dim), dtype=complex)
    for i in range(3):
        for j in range(3):
            src = i*3 + j
            if p_idx == 0:
                tgt_i, tgt_j = i + sign, j
            else:
                tgt_i, tgt_j = i, j + sign
            if 0 <= tgt_i < 3 and 0 <= tgt_j < 3:
                tgt = tgt_i*3 + tgt_j
                mat[tgt, src] = 1.0
    return qt.Qobj(mat)

L2p, L2m = L_p(0, +1), L_p(0, -1)
L3p, L3m = L_p(1, +1), L_p(1, -1)

# Example: assemble Lindblad operators for one baseline
gamma = 0.5   # dephasing strength
kappa = 1.0   # repair strength

# Dephasing collapse operators (diagonal in Q2, Q3 basis)
c_ops_dephase = [np.sqrt(gamma) * Q2, np.sqrt(gamma) * Q3]

# Repair collapse operators (biased choice of L2p, L3p, etc.)
c_ops_repair = [np.sqrt(kappa) * L2p, np.sqrt(kappa) * L3p]

# Full collapse operator list for the multiplicity baseline
c_ops_mult = c_ops_dephase + c_ops_repair

# Hamiltonian (placeholder; can be taken as 0 or a simple diagonal)
H = qt.Qobj(np.zeros((dim, dim), dtype=complex))

# Initial state in the engram band
rho0 = engram_states[0] * engram_states[0].dag()

# Time evolution
tlist = np.linspace(0, 10.0, 200)
result = qt.mesolve(H, rho0, tlist, c_ops_mult, [])

# Compute retention in H_E at final time
rho_final = result.states[-1]
retention = (Pi_E * rho_final * Pi_E).tr()
print("Engram-band retention:", np.real_if_close(retention))
\end{verbatim}

The two comparison baselines modify only the \texttt{c\_ops} set and the
interpretation of the basis, keeping dimension, adjacency, and resource
budget fixed.

\subsection{Falsification Criterion}

A sharp falsification criterion is:

\begin{quote}
If under matched Hilbert-space dimension, leakage graph, and control budget
(number and strength of collapse channels), the multiplicity code does not
achieve strictly higher engram-band retention (or logical fidelity) than a
generic lattice-label code and a stabilizer-like code, then the prime-graded
structure has not yet demonstrated a unique physical advantage.
\end{quote}

Conversely, a robust and repeatable advantage under these constraints
strengthens the claim that multiplicity grading is a fundamental symmetry
class rather than a mere relabeling.

\section{Discussion and Future Work}

\subsection{Theoretical Extensions}

Natural extensions include:
\begin{itemize}[leftmargin=*]
  \item Generalization to larger prime sets and higher-dimensional bands
        $\mathcal B$ with richer combinatorial structure.
  \item Rigorous proofs of distance and threshold properties under specific
        noise models, analogous to stabilizer-code threshold theorems.
  \item Formalization in proof assistants (e.g.\ Lean) to align with existing
        formalization efforts in multiplicity-based spectral problems.
\end{itemize}

\subsection{Experimental Directions}

Beyond simulation, minimal quantum hardware tests could be performed on small
superconducting or ion-trap devices, where:
\begin{itemize}[leftmargin=*]
  \item two or more commuting observables play the role of low-prime charges;
  \item leakage channels are engineered by controlled noise or sequences of
        gates;
  \item multiplicity-biased repair is implemented by conditional pulses or
        measurement-based feedback.
\end{itemize}

Parallel EEG/MEG analysis using prime-locked pipelines provides an
independent, macroscopic test bed for multiplicity-graded stability,
complementing the microscopic quantum experiments.

\section*{Conclusion}

We have formulated a multiplicity engram code as a prime-graded subsystem
code, connecting prime-encoded memory states, tensor-network structure, and
multiplicity feedback into a single operator-algebraic framework. The
definition, dynamics, distance, and non-reducibility statements are all
consistent with existing Multiplicity Theory work and the \emph{Quantum
Memory Engrams} paper, and they lead directly to concrete simulation and
hardware experiments that can corroborate or falsify the claimed advantage
of multiplicity-graded codes.

\nocite{*}
\bibliographystyle{plain}
\bibliography{references}

\end{document}